\documentclass[11pt]{article}
\usepackage[T1]{fontenc}
\usepackage[utf8]{inputenc}
\usepackage{amsmath,amssymb,amsthm,booktabs,longtable,array}
\usepackage[margin=1in]{geometry}
\usepackage{hyperref}
\usepackage{xurl}
\newtheorem{theorem}{Theorem}[section]
\newtheorem{lemma}[theorem]{Lemma}
\newtheorem{proposition}[theorem]{Proposition}
\newtheorem{definition}[theorem]{Definition}
\newtheorem{certificate}[theorem]{Finite certificate lemma}
\newcommand{\Q}{\mathbb Q}
\newcommand{\HH}{\mathbb H}
\newcommand{\wt}{\operatorname{wt}}
\newcommand{\rk}{\operatorname{rk}}
\newcommand{\diag}{\operatorname{diag}}
\title{A computer-assisted dimension bound for compact arithmetic\linebreak hyperbolic Coxeter polytopes}
\author{Coxeter Polytopes Frontier Project}
\date{Proof translation, 18 October 2026}
\begin{document}
\maketitle
\begin{abstract}
We give an exposition of the completed internal proof that a compact arithmetic
hyperbolic Coxeter polytope has dimension at most 21. The proof combines
published high-dimensional field restrictions with complete local arithmetic
reductions and exact rational angle weights. The difficult golden-field step
excludes dimension 22 by a vertex bound $21977/500<44$; the strict lower bound 44
comes from the $h$-vector of a simple compact polytope. We explain why the
local classification covers arbitrary supporting walls, why singular
intermediate states must remain payment obligations, and why the normalized
cross-component weights admit exact forest maximization. This document
translates the existing proof and its certificate interfaces; it does not
report a new geometric or arithmetic census. The result is internally checked,
not externally refereed or fully formalized in a proof assistant. No
compact dimension 21 example or compact-only sharpness is asserted.
\end{abstract}

\section{Statement, scope, and proof architecture}
A Coxeter polytope means a finite-sided convex hyperbolic polytope whose
facet dihedral angles are integral submultiples of $\pi$. Arithmeticity
refers to the lattice generated by its facet reflections and its actual
Vinberg ground field, not to a field merely containing displayed entries.
\begin{theorem}\label{thm:main}
Every compact arithmetic hyperbolic Coxeter polytope has dimension at most 21.
\end{theorem}
The internal proof of this theorem is complete: the bound-21 composition,
its independent scope checker, and the geometric audit identify no unresolved
mathematical step. Historical immutable reports use ``proof candidate'' to
record the absence of external validation; this is not an unfilled step in
the internally completed argument. No priority claim is made here.

The proof has three kinds of inputs, which must not be confused.
\begin{enumerate}
\item Published geometry and arithmetic: compact simplicity and face
correspondence, finite, affine and Lann\'er classifications, arithmetic cyclic-product
conditions, the compact dimension ceiling 29, the high-dimensional ground-field
list, and the nonnegative symmetric $h$-vector theorem.
\item Internal mathematical arguments: projection to actual faces,
complete-star rank, first-touch reduction, all-size attachment closure,
whole-component weight normalization, and the strict counting contradiction.
\item Finite exact certificates: every member of a proved finite domain is
checked for its necessary arithmetic conditions, literal reduction maps,
face payments, and finite-forest upper bounds. Floating optimization proposes
weights but is not used to certify them.
\end{enumerate}
The detailed assumption/reproduction map is
\path{docs/compact_arithmetic21_proof_map.md}. The final composition is
\path{scripts/fresh_arithmetic_dimension21_composition_20261016.py}; the
independent final scope check is
\path{scripts/fresh_arithmetic_dimension21_scope_check_20261016.py}.
Neither small program by itself establishes the geometric premises below.

\section{The strict weighted-face contradiction}
Let $P$ be a simple compact $n$-polytope, and write $f_i$ for its numbers of
$i$-faces. A planar angle is a pair $(v,F)$ where $v$ is a vertex of a
2-face $F$. At $v$ it corresponds to an unordered pair of the $n$ incident
facets, equivalently a marked pair in its finite Coxeter diagram.
\begin{lemma}\label{lem:count}
Suppose nonnegative angle weights satisfy
\[
 \sum_{v\in F}\wt(v,F)\ge 5-k
 \quad\text{for every $k$-gonal 2-face $F$}.
\]
Then the average total vertex weight is strictly greater than
\[
 \beta(n)=
 \begin{cases} n(n-6)/8,&n\text{ even},\\
 (n-1)(n-5)/8,&n\text{ odd}.
 \end{cases}
\]
In particular, a uniform vertex upper bound at most $\beta(n)$ is impossible.
\end{lemma}
\begin{proof}
Double counting angles gives
\[
 \sum_v\wt(v)\ge 5f_2-\binom n 2 f_0.
\]
For the nonnegative symmetric $h$-vector of the simplicial dual,
\[
 f_0=\sum_{i=0}^n h_i,\qquad
 f_2=\sum_{i=0}^n\binom{n-i}{2}h_i.
\]
Pair $i$ and $n-i$. The minimum paired coefficient is
$n(n-2)/8$ for even $n$ and $(n-1)^2/8$ for odd $n$.
The endpoint coefficients are strictly larger, and $h_0=h_n=1$.
Consequently $f_2/f_0$ is strictly larger than that minimum. Substitution
proves the claimed formula. Strictness does not require every middle
$h_i$ to be positive.
\end{proof}
It suffices to pay 2 on every triangle and 1 on every quadrilateral;
nonnegativity handles every $k\ge 5$. We never assume a face is itself a
Coxeter polytope: its angles are evaluated from the ambient finite vertex
forest.

\section{Actual faces and complete supporting stars}
Use the Lorentz convention $(n,1)$, inward chamber inequalities
$\langle x,r_i\rangle\le 0$, and spacelike outward roots with pairwise
nonpositive products. A finite corner has positive-definite Gram matrix.
\begin{lemma}\label{lem:face}
An elliptic principal root set determines an actual face of the compact
chamber. A positive affine principal set is incompatible with compactness.
\end{lemma}
\begin{proof}
For the elliptic set let $B$ be its positive-definite Gram matrix and $h$
a strict interior timelike vector. Since $B$ is a positive-definite
$Z$-matrix, $B^{-1}$ is entrywise nonnegative. Put
$c=B^{-1}(\langle h,r_i\rangle)_i\le 0$ and
$h_0=h-\sum_i c_i r_i$. It is orthogonal to the chosen roots,
\[
 h_0^2=h^2-(\langle h,r_i\rangle)^TB^{-1}
 (\langle h,r_i\rangle)<0,
\]
and has strictly negative product with every other chamber root. This
supplies a relative-interior face point. For an affine connected principal
set, its positive kernel produces a nonzero null vector: it cannot be the
zero vector because pairing the positive relation with $h$ is negative.
Its products with all other roots are nonpositive. It gives an ideal point
of the chamber, contrary to compactness.
\end{proof}

For a quadrilateral 2-face, let $W$ be its full $n-2$ supporting roots and
let $b_0,b_1,b_2,b_3$ be its cyclic boundary roots. Every
$W+\{b_i,b_{i+1}\}$ is spherical. Each opposite extension
$W+\{b_0,b_2\}$ or $W+\{b_1,b_3\}$ is indefinite by
Lemma~\ref{lem:face}; an elliptic opposite extension would make opposite
facets meet on this face. Select a minimal indefinite subdiagram in each.
It must contain both relevant black roots because deleting either leaves
a principal subset of a finite corner. These are arithmetic Lann\'er
pieces; the order bound 5 is a published input.

The union of the two pieces is only an initial core. Shared white roots,
all black/white identifications, and every cross contact are retained.
Every unspecified pair belongs to some finite corner and therefore has a
finite label. Opposite dotted contacts remain arithmetic parameters.

\begin{lemma}\label{lem:rank}
A $k$-root core with a positive-definite corner of order $k-2$ and a
negative direction has nullity at most 1. The complete $n+2$-root
quadrilateral star has nullity exactly 1.
\end{lemma}
\begin{proof}
The positive index is at least $k-2$, and the physical ambient form has
negative index 1. This proves the first assertion. A complete star has
$k=n+2$ and lies in Lorentz dimension $n+1$, so its Gram is singular;
the first assertion makes its nullity exactly 1.
\end{proof}

If a core is singular and has a negative direction, its Gram kernel is an
actual relation among its normals: a Lorentz subspace containing a timelike
vector is nondegenerate. Every later contact row must annihilate that
kernel. A nonnegative root relation is impossible by the strict interior
point. These are exact geometric rejection tests, not numerical rank
heuristics.

Every additional white component touching a core has a first touching
root, followed by an ordering along its connected white edges. Enumerating
all such one-white contacts is therefore enough for a complete attachment
queue. If no unpaid contact survives, all remaining roots are orthogonal
spectators. Marked connected components in the core are then entire in
the full finite corner. A nonsingular core plus only positive-definite
orthogonal spectators cannot acquire the nullity required by
Lemma~\ref{lem:rank}.

A singular state may stop with only orthogonal spectators \emph{even if it
also has possible touching extensions}. Thus every singular intermediate
state must be retained as a separate payment obligation. Keeping only
terminal sealed states would be an incorrect completeness argument.

\section{The two fields with an all-rank bound 20}
The arithmetic condition used throughout is integral cyclic products in
the actual ground field and positive-semidefinite conjugate Grams at every
nonphysical embedding. Finite contacts and Lann\'er pieces are classified
using these conditions, rather than a chosen label cutoff.

\subsection{The field \texorpdfstring{$\Q(\sqrt 2)$}{Q(sqrt2)}}
The complete finite alphabet is $2,3,4,6,8$: the two-cycle
$2+2\cos(2\pi/m)$ must belong to the field. The Lann\'er list consists of
an arbitrary permitted dotted pair, nine triangles, and the two chordless
cycles $C_4,C_5$ with one 4-edge and all remaining edges 3. The triangle
labels are
\[
\begin{gathered}
(2,3,8),(2,4,8),(2,6,8),(2,8,8),(3,3,4),\\
(3,8,8),(4,4,4),(4,6,6),(4,8,8).
\end{gathered}
\]
The field catalogue proves completeness through cyclotomic degree,
squareclass parity and exact conjugate determinants.

Assign a base unit to every marked pair at graph distance at most 2.
For the unpaid case, all marked minimal unions are exhausted: 51 marked
pieces, 6,010 shared-white identifications, 115 literal unions and 52
marked classes. Forty classes have two negative directions and one fails
cyclic field membership. Nine finite-label cores and two dotted families
remain. Complete one-white domains leave seven singular finite cores
sealed; a further nonsingular core can only extend to a sealed core.
The last finite core admits only short-arm path extensions of $D_5$ to
$E_6,E_7,E_8$, or a first affine obstruction. If its two black endpoints
are not connected through new whites, elimination gives reduced unit
squares
\[
 p,q\in\{1,3/4,2/3,5/8\},
\]
and complete-star singularity requires
\[
 (p+q-4pq)+(1-p-q)\sqrt 2=0.
\]
This would force $p+q=1$, contradicting $p+q\ge 5/4$.
Connecting paths give literal sealed-core maps, not additional unbounded
families.

For the dotted cores, every outside contact is at one of the two dotted
endpoints. Eliminating a white component joining them leaves
$0<p,q\le 3/4$ and increases the effective norm-two dotted magnitude to
$r_{\rm eff}>2$. In the single-contact case the rank identity gives
\[
 r_{\rm eff}^2=q(4p-1+\sqrt 2)
 \le\tfrac 34(2+\sqrt 2)<4.
\]
In the double-contact case it gives
\[
 r_{\rm eff}\le\tfrac 12+\sqrt 2<2.
\]
Both are contradictions. The independent pendant classification covers
all remaining disconnected tails. Thus no arbitrary white component has
been silently replaced by a bounded path search.

Add whole-component bonuses to the base rule:
\begin{center}
\begin{tabular}{lll}\toprule
Component&Marked pairs&Total bonus\\\midrule
$F_4$&two path endpoints&1\\
$D_5$&long tip with either short tip&1\\
$E_6$&near/far nodes on the two equal arms&1\\
$E_7,E_8$&specified unequal-arm near/far pair&$1/4$ each\\\bottomrule
\end{tabular}
\end{center}
The named unequal-arm pair is the near node of the longest arm and the
length-two-arm tip. The $D_5,E_6$ pairs each receive halves.
The sealed-corner tables give every exceptional quadrilateral at least 1.
Triangles receive at least 2: a three-black minimal Lann\'er triangle has
two nonright edges, while a four- or five-cycle has at least two black
cyclic gaps of length at most 2. These paths remain in the corresponding
finite corners under additional white extensions.

The distance-two count on a rank-$r$ path is $2r-3$, and on a $D/E$ tree
it is $2r-2$. The only deficits below 2 from $2r$ after bonuses are
$1,1,7/4,7/4$ on $D_5,E_6,E_7,E_8$. For total rank $n\ge 9$, a
single component has deficit at least 2, or at least two components have
combined deficit at least 2. Every vertex therefore has weight at most
$2n-2$.

\subsection{The cubic field \texorpdfstring{$K_7=\Q(\cos(2\pi/7))$}{K7}}
Here the arithmetic Lann\'er classification leaves orders 2 and 3. A
triangle has a distinguished 7- or 14-edge; without such an edge its
entries and determinant field cannot generate $K_7$. A high edge in a
finite diagram is isolated. Complete minimal-union finite-corner and
short-path conditions leave nine supports. The arbitrary-white argument
projects off the two opposite positive black blocks: white neighbors
have negative residual squares; a 14-edge or a dense contact would give
two negative directions. All contacts between the two surviving white
neighbor sets are nonright, so the white forest forces one set to be a
singleton and the other to have size 1 or 2. Off-core components attach
only once and have rational Schur diagonal corrections.

In the six-node case, full-star singularity would give a nonzero rational
quadratic in the degree-three number $c^2$, which is impossible. In the
seven-node case write the remaining white matrix as
\[
 H=\begin{pmatrix}p&-\delta&-\varepsilon\\
 -\delta&q&0\\-\varepsilon&0&r\end{pmatrix}>0,
 \qquad L=\delta\varepsilon\beta\gamma.
\]
Cycle-field membership puts $L$ in $K_7\cap\Q(\sqrt 2)=\Q$.
The coefficient of the linear term in the reduced cubic determinant is
$(\det H-8L)/8$. Consequently
\[
 8L=\det H\le 1-\delta^2-\varepsilon^2\le\tfrac 12\le 8L.
\]
Equality removes every external branch and forces the displayed contacts
to be simple. The other coefficient forces the remaining contact simple.
The two exceptional paying corners are entire $D_5$ factors.

Use the distance-two base, one additional unit on isolated
$I_2(7),I_2(14)$, and halves on the two long-tip/short-tip pairs of whole
$D_5$. Triangles receive 2 at a high-edge corner. Quadrilaterals receive 1
from a short corner or the two $D_5$ halves. The same component-deficit
argument gives vertex weight at most $2n-2$ for $n\ge 6$.

\begin{proposition}\label{prop:twofields}
For either $K=\Q(\sqrt 2)$ or $K=K_7$, a compact arithmetic Coxeter
polytope has dimension at most 20.
\end{proposition}
\begin{proof}
The finite certificates and all-size arguments just described establish
the hypotheses of Lemma~\ref{lem:count}. At $n=21$ its strict lower bound
is 40, equal to the upper $2n-2$. More generally, for $n=21+2u$ or
$n=22+2u$, subtraction of $2n-2$ from $\beta(n)$ gives respectively
$u^2/2+5u$ and $u^2/2+11u/2+2$. Both are nonnegative for every integer
$u\ge 0$. Strictness excludes every $n\ge 21$.
\end{proof}
The finite numerical identities in this section are certified by the
independent square-root-field quadrilateral and weight reviewers and the
cubic-seven global reviewer listed in the proof map. Their all-size
geometric proofs are essential; the tables alone would not suffice.

\section{Complete golden-field reduction}
Put $\phi=(1+\sqrt 5)/2$ and $K=\Q(\phi)$.
The finite alphabet is $2,3,4,5,6,10$. The exact arithmetic Lann\'er
catalogue has 12 triangles, six order-four types, four order-five types,
and the dotted parameter
\[
 t\in\mathbb Z[\phi],\qquad t>4,\quad 0<t'<4.
\]
The prime denotes the other real embedding. The fixed baseline is
\[
 b(C;i,j)=\mathbf 1_{d_C(i,j)\le 2}
 +\mathbf 1_{C\cong H_4}\,
    \mathbf 1_{\{d_C(i,j)=3\}\ \text{or}\ \{\text{$ij$ is the 5-edge}\}}.
\]
Here $C$ is an \emph{entire} connected component. Short paths persist in
finite extensions, and $H_4$ has no larger connected spherical extension.
The 76 minimal triangle markings give baseline at least 2.

For an unpaid quadrilateral, the no-distance-two minimal-union reduction
has 5,440 literal supports and 1,180 marked classes. Arithmetic and
signature filters leave 362 finite-label cores and 34 dotted families.
After persistent payments and affine exclusions, the obligations are:
\begin{center}
\begin{tabular}{lr}\toprule
Reduction&Number\\\midrule
Original sealed singular rows&112\\
Singular first attachments&147\\
Marked singular frontier rows&51\\
Further closure rows&17\\
Singular dotted-queue stopping rows&264\\
Singular finite-queue stopping rows&2,070\\\midrule
Total zero-touch obligations&2,661\\\bottomrule
\end{tabular}
\end{center}
Every listed singular row is retained even if it also admits touching
extensions. Nonsingular orthogonally sealed rows are rank-excluded.

The finite queue has 5,172 states and 51,599 complete one-white profiles.
The dotted queue has 438 states, including 423 below its geometric boundary,
and 5,146 next-white profiles. Each retained child has a literal full-label
bijection; marked black pairs and exact dotted parameters are preserved.
The maximum 31 nodes is justified by $n+2\le 31$ after the published
compact ceiling 29, not by a resource cutoff.

For the 24 residual dotted families, deleting the two dotted endpoints
leaves a nonsingular Lorentz block $B$. The endpoint Schur complement
must be positive semidefinite in any ambient Lorentz extension. Writing
it as $\left(\begin{smallmatrix}a&-x-c\\-x-c&b\end{smallmatrix}\right)$
gives $x^2\le 2ab+2c^2$. Arithmeticity bounds the corresponding integral
two-cycle $t=A+B\phi$ simultaneously at the two embeddings. This gives
a finite integer box, containing 532 tested candidates and 32 admissible
parameter seeds. Hence the dotted queue is exact and finite; no distance
sampling replaces a continuous parameter argument.

The first-touch argument of Section 3 proves that every actual unpaid
quadrilateral reaches a retained stopping row or a persistent baseline
payment. Remaining orthogonal white components are arbitrary spectators.
This is the geometric content of complete coverage, beyond a statement
that the stored queue has no pending jobs.

\section{Exact whole-component weights at rank 22}
Let $a,b,c$ count entire $A_1,I_2(5),H_3$ components of a finite vertex
forest. The rational coefficient table is specified by
\path{results/fresh_sqrt5_dimension22_h3_cross_20261016_cap2.json}:
its \texttt{variable\_order} indexes \texttt{certificate.exact\_primal}.
The terminal coordinate records the vertex upper; the preceding 30,169
coordinates are nonnegative rational bonuses. There are 232 connected
marked-pair orbits,178 rooted $A_1$ destinations,97 rooted $I_2(5)$
destinations and six rooted $H_3$ destinations. Orbit graphs and marks,
not opaque string keys alone, define the table.

The coefficient context is $(a,\min(b,2),\min(c,1))$. Add a connected
bonus to a baseline-zero marked pair in the same entire component. For
pairs in distinct components use the following disjoint rules:
\begin{center}
\begin{tabular}{lll}\toprule
Source&Excluded destination types&Payment denominator\\\midrule
$A_1$&$A_1$&$a$\\
$I_2(5)$&$A_1,I_2(5)$&$2b$\\
$H_3$&$A_1,I_2(5),H_3$&$3c$\\\bottomrule
\end{tabular}
\end{center}
The numerator is the rooted destination coefficient in that context.
All omitted coefficients are zero. In particular an $A_1$--$H_3$ pair
belongs only to the first rule, and an $I_2(5)$--$H_3$ pair only to the
second. Denominators use \emph{actual} counts, not capped counts.

\begin{lemma}\label{lem:additive}
At fixed exact $(a,b,c)$ the total vertex charge is additive over entire
finite components.
\end{lemma}
\begin{proof}
For any eligible destination root there are respectively $a,2b,3c$
source roots. Their contributions sum to the numerator once. Disjointness
prevents double counting. Connected bonuses and the baseline are already
componentwise. A hub component may still receive incoming payments from
an earlier hub family; these are included in its own additive cost.
\end{proof}

Let $C_T(a,b,c)$ denote the resulting exact cost of type $T$. Fix the
three hub counts and put $r=22-a-2b-3c$. The residual forest contains none
of the three hub types. With $D(0)=0$ and unattainable ranks assigned
$-\infty$, use
\[
 D(s)=\max_{T\notin\{A_1,I_2(5),H_3\},\,\rk T\le s}
       \{D(s-\rk T)+C_T(a,b,c)\}.
\]
This recurrence includes every multiset of finite components by removing
one final component; conversely every transition appends a permitted
component. Add the costs of the fixed hubs at the end. Rank 1 is the sole
unattainable residual rank, since $A_2,A_3$ realize every integer at least 2.

For a $k$-node stopping core at ambient rank 22, spectator rank is $24-k$.
Enumerate every common external triple $(e,f,g)$ with
\[
 e+2f+3g\le 24-k,\qquad 24-k-e-2f-3g\ne 1.
\]
The same external triple is added to each corner's own internal counts
before the coefficient context and denominator are evaluated. Choosing
spectators independently at the four corners would not prove a face
inequality.

\begin{certificate}\label{cert:golden}
The specified rank 22 rational rule has the following exact properties.
\begin{enumerate}
\item All 76 triangle baseline totals are at least 2.
\item All 2,661 stopping identities are represented. Precisely 1,634 fit
rank 22; their 91,726 common spectator triples have nonnegative corner
weights and quadrilateral total at least
\[
\frac{560003}{560000}>1.
\]
\item All 71 finite component types and 410 exact hub-count contexts have
vertex total at most $21977/500$.
\end{enumerate}
\end{certificate}
The independent reviewer reconstructs finite graphs and rooted/marked
orbits, tests every literal label-preserving map, and runs a type-first
unbounded-coin recurrence in reverse component order, distinct from the
producer's rank-first recurrence. The complete coverage program
reevaluates all older stopping branches under this new rule, rather than
assuming an old numerical payment transfers. The certificate lemma is a
finite exact computation on an explicitly complete domain, with the
published classifications and Section 3 supplying its geometric meaning.

At $n=22$, Lemma~\ref{lem:count} gives strict average weight greater
than 44. Certificate~\ref{cert:golden} gives at most
$21977/500=43.954$. This contradiction excludes dimension 22.
The already completed dimension-specific rules exclude all larger
remaining dimensions:
\begin{center}
\begin{tabular}{rrr}\toprule
$n$&Exact vertex upper&Strict lower threshold $\beta(n)$\\\midrule
22&$21977/500$&44\\
23&$485211/10000$&$99/2$\\
24&$267527/5000$&54\\
25&59&60\\
26&$31017/500$&65\\
27&$67009/1000$&$143/2$\\
28&$230/3$&77\\
29&84&84\\\bottomrule
\end{tabular}
\end{center}
Each row has its own fixed rational rule and complete current-rule face
payments on the same exhaustive arithmetic reduction. The equality at 29
is contradicted by the strict lower bound. One does not extrapolate the
rank 22 coefficients to higher ranks.
\begin{proposition}\label{prop:golden}
Compact arithmetic Coxeter polytopes over $\Q(\sqrt 5)$ have dimension
at most 21.
\end{proposition}
\begin{proof}
The table excludes 22 through 29. The published compact ceiling excludes
higher dimensions.
\end{proof}

\section{Field composition and the noncompact corollary}
\begin{proof}[Proof of Theorem~\ref{thm:main}]
Suppose $n\ge 22$. Vinberg's compact ceiling gives $n\le 29$.
The published high-dimensional field restriction~\cite{Vinberg,Nikulin} leaves
$\Q,\Q(\sqrt 2),\Q(\sqrt 5),K_7$. The rational field is impossible
in this compact range: arithmetic reflection lattices are of simplest
type, and a rational indefinite form in at least five variables is
isotropic by Meyer, so its arithmetic quotient is noncompact.
Proposition~\ref{prop:twofields} excludes the two bound 20 fields, and
Proposition~\ref{prop:golden} excludes the golden field. This covers all
$8\cdot 3=24$ compact field/dimension branches, using the field restriction
only in its stated range $n\ge 22$.
\end{proof}

Combining this with Esselmann's published noncompact arithmetic bound~\cite{Esselmann}
$n\le 21$ gives the same bound for all finite-volume arithmetic reflection
groups. Borcherds' noncompact dimension 21 example~\cite{Belolipetsky} makes that combined
bound sharp. It does not make the compact-only bound sharp. No result
here concerns arbitrary nonarithmetic compact chambers beyond the
separately known ceiling 29.

\section{Verification boundaries and further work}
The complete proof is computer-assisted, not a mere manifest comparison.
Its mathematical parts are the reductions and counting implications above;
its finite parts are exact field/inertia, graph-domain, map, weight and
forest checks. Hashes identify which finite objects were checked. They do
not prove geometric completeness. A separate final scope checker matches
all 2,661 obligation identities,91,726 common triples,410 contexts and 24
field/dimension branches to the originating catalogues.

This translation does not rerun the large classification queues. The
reproduction map separates a quick immutable binding/scope replay from
full independent finite-check replays. Published inputs may be axioms in
a future proof-assistant layer. Internal certificate lemmas must instead
be proved, reflected into checked computation, or explicitly reported as
remaining formalization assumptions. A Lean theorem importing the entire
field bounds as axioms would certify only the final composition, not this
whole computer-assisted proof.

Compact dimension 21 remains an existence/nonexistence problem; later
field-specific work reduces its possible ground field to $\Q(\sqrt 5)$.
The current rank 13 incidence model has a verified dual lower bound above
its target 12, so additional valid incidence information is required;
reoptimizing the same matrix is insufficient. In the general compact
lane, excluding dimension 29 via this weighting framework requires a
vertex upper at most $\beta(29)=84$ instead of the classical $3n=87$.
A concrete next step is to classify the maximizing rank 29 forests under
the current geometric rule and seek whole-component exchanges whose
face proofs preserve every triangle and quadrilateral payment. These
are research directions, not improvements asserted by this manuscript.

\begin{thebibliography}{9}
\bibitem{Vinberg} E. B. Vinberg, \emph{Absence of crystallographic groups of
reflections in Lobachevskii spaces of large dimension}, Transactions of the
Moscow Mathematical Society 47 (1984),75--112; original Russian text,
\url{https://www.mathnet.ru/php/getFT.phtml?jrnid=mmo&option_lang=rus&paperid=444&what=fullt}.
\bibitem{Nikulin} V. V. Nikulin, \emph{On ground fields of arithmetic
hyperbolic reflection groups}, especially Introduction and Sections 2,3,
\url{https://arxiv.org/abs/0708.3991}.
\bibitem{Belolipetsky} M. Belolipetsky, \emph{Arithmetic hyperbolic reflection
groups}, especially Lemma 2.1 and Sections 1,7,
\url{https://arxiv.org/abs/1506.03111}.
\bibitem{Esselmann} F. Esselmann, \emph{\"Uber die maximale Dimension von
Lorentz-Gittern mit coendlicher Spiegelungsgruppe}, Journal of Number
Theory61 (1996),103--144, \url{https://doi.org/10.1006/jnth.1996.0141}.
\bibitem{Internal} Exact internal sources and independent reviews, fixed
16 October 2026: the path-indexed inventory in
\path{docs/compact_arithmetic21_proof_map.md}. Historical status wording is
preserved in the frozen files; this translation supplies no new priority,
external-referee, or complete proof-assistant claim.
\end{thebibliography}
\end{document}
